%/*% pdflatex -file-line-error -interaction nonstopmode %

\documentclass{article}
\usepackage{graphicx}
\usepackage{amsmath}
\usepackage{amsfonts}
\usepackage{bookmark}
\usepackage{booktabs}
\usepackage{exercises}
\graphicspath{ {./images/test} }
\usepackage{indentfirst}
\usepackage{amsthm}
\usepackage{yhmath}
\usepackage{placeins}
\newcommand*{\oldneg}{\mathord{\sim}}
\newcommand\showdiv[1]{\overline{\smash{)}#1}}
\newcommand\Mydiv[1]{%
\smash{\raise.3ex\hbox{$\big)$}}\mkern-8mu
        \overline{\enspace\strut#1}}
\setlength{\tabcolsep}{5pt}
\theoremstyle{definition}
\newtheorem{definition}{Definition}

\setlength{\parindent}{.5in}
\setlength{\parskip}{12pt}
\renewcommand{\arraystretch}{2}

\title{Quiz 2 Answers (pg. 197)}
\begin{document}
\maketitle

\section*{1}
\[ f(x)+g(x)=4x^2-x+2x^2=6x^2-x \]
Because the domains of both $f$ and $g$ are all real numbers, the domain of $f(x)+g(x)$ is all real numbers.

\section*{2}
\[ g(x)-f(x)=2x^2-(4x^2-x)=-2x^2+x \]
Like the first question, the domain of $g(x)-f(x)$ is all real numbers.

\section*{3}
\[ f(x)\times g(x)=(4x^2-x)(2x^2)=8x^4-2x^3 \]
Again, the domain of $f(x)\times g(x)$ is all real numbers.

\section*{4}
\[ \frac{f(x)}{g(x)}=\frac{4x^2-x}{2x^2}=\frac{4x-1}{2x} \]
The domain of $\frac{f(x)}{g(x)}$ is all real numbers other than 0, as we cannot divide by $2(0)=0$.

\section*{5}
\[ f(g(x))=f(2x^2)=4(2x^2)^2-2x^2=4(4x^4)-2x^2=16x^4-2x^2 \]
Again, the domain of $f(g(x))$ is all real numbers.

\section*{6}
\[ g(f(x))=g(4x^2-x)=2(4x^2-x)^2=2(16x^4-8x^3+x^2)=32x^4-16x^3+2x^2 \]
Which also has a domain of all real numbers.

\section*{7}
\[ f(f(x))=f(4x^2-x)=4(4x^2-x)^2-(4x^2-x)=4(16x^4-8x^3+x^2)-4x^2+x \]
\[ =64x^4-32x^3+4x^2-4x^2+x=64x^4-32x^3+x \]
which also has a domain of all real numbers.

\section*{8}
\[ g(g(x))=g(2x^2)=2(2x^2)^2=2(4x^4)=8x^4 \]
which also has a domain of all real numbers.

\section*{9}
We must show that $f(g(x))=x$ and $g(f(x))=x$.
\[ f(g(x))=f(x+9)=x+9-9=x \]
\[ g(f(x))=g(x-9)=x-9+9=x \]

\section*{10}
We must show that $f(g(x))=x$ and $g(f(x))=x$.
\[ f(g(x))=f(\sqrt[3]{\frac{x}{5}})=5(\sqrt[3]{\frac{x}{5}})^3=5\frac{x}{5}=x \]
\[ g(f(x))=g(5x^3)=\sqrt[3]{\frac{5x^3}{5}}=\sqrt[3]{x^3}=x \]

\section*{11}
We must show that $f(g(x))=x$ and $g(f(x))=x$.
\[ f(g(x))=f(-\frac{2}{3}x+\frac{1}{6})=-\frac{3}{2}(-\frac{2}{3}x+\frac{1}{6})+\frac{1}{4}=x-\frac{1}{4}+\frac{1}{4}=x \]
\[ g(f(x))=g(-\frac{3}{2}x+\frac{1}{4})=-\frac{2}{3}(-\frac{3}{2}x+\frac{1}{4})+\frac{1}{6}=x-\frac{1}{6}+\frac{1}{6}=x \]

\section*{12}
We must show that $f(g(x))=x$ and $g(f(x))=x$. 
\[ f(g(x))=f((\frac{x-1}{6})^\frac{1}{2})=6((\frac{x-1}{6})^\frac{1}{2})^2+1=6(\frac{x-1}{6})+1=x-1+1=x \]
\[ g(f(x))=g(6x^2+1)=(\frac{6x^2+1-1}{6})^\frac{1}{2}=(x^2)^\frac{1}{2}=x \]

\section*{13}
\[ y=-\frac{1}{3}x+5 \textrm{ (Given function)} \]
\[ x=-\frac{1}{3}y+5 \textrm{ (Switch variables)} \]
\[ x-5=-\frac{1}{3}y \]
\[ -3x+15=y \]

\section*{14}
\[ y=x^2-16,x\ge 0 \]
\[ x=y^2-16 \]
\[ x+16=y^2 \]
\[ \pm\sqrt{x+16}=y \]
Because the domain of the initial function is $x\ge 0$, the inverse of the function is $y=+\sqrt{x+16}$.

\section*{15}
\[ y=-\frac{2}{9}x^5 \]
\[ x=-\frac{2}{9}y^5 \]
\[ -\frac{9}{2}x=y^5 \]
\[ \sqrt[5]{-\frac{9}{2}x}=y \]

\section*{16}
\[ y=5x+12 \]
\[ x=5y+12 \]
\[ x-12=5y \]
\[ \frac{x-12}{5}=y \]

\section*{17}
\[ y=-3x^3-4 \]
\[ x=-3y^3-4 \]
\[ x+4=-3y^3 \]
\[ -\frac{x+4}{3}=y^3 \]
\[ \sqrt[3]{-\frac{x+4}{3}}=y \]

\section*{18}
\[ y=9x^4-49,x\le 0 \]
\[ x=9y^4-49 \]
\[ x+49=9y^4 \]
\[ \frac{x+49}{9}=y^4 \]
\[ \pm\sqrt[4]{\frac{x+49}{9}}=y \]
Because $x\le 0$ in the original function, the inverse of the function is $y=-\sqrt[4]{\frac{x+49}{9}}$.

\section*{19}
\[ C(g(d))=C(0.02d)=2.15(0.02d)=0.043d \]
\[ C(g(400))=C(0.02(400))=C(8)=2.15(8)=17.20 \]
$C(g(400))$ represents the cost in dollars for the amount of gallons of gasoline needed to travel a distance of 400 miles.

\end{document}
